\documentclass[11pt]{article}
\usepackage[a4paper,margin=2.3cm]{geometry}
\usepackage{xcolor,amsmath,amssymb}
\usepackage[hidelinks]{hyperref}
\newcommand{\verify}[1]{\textcolor{red}{[VERIFY: #1]}}
\newcommand{\rc}[1]{\par\medskip\noindent\textit{#1}\par\smallskip}
\newcommand{\resp}{\noindent\textbf{Response.} }
\setlength{\parindent}{0pt}
\title{Response to the reviewer\\[4pt]\large What Does EEG-to-fMRI Synthesis Add to EEG-Based Diagnosis? A Controlled Benchmark in Psychosis, Schizophrenia and Parkinson's Disease}
\date{}
\begin{document}
\maketitle

We thank the reviewer for a careful and constructive review. Every point has been addressed with new analyses or changes to the text; the main changes are summarised first, followed by point-by-point responses. Section, table and figure numbers refer to the revised manuscript. All analyses use the same subjects, splits, classifier and nested cross-validation as before, and all code is released.

\section*{Summary of the main changes}
\begin{enumerate}
\item Two EEG-only baselines matched in information to the synthesized view (the temporal mean and SD of the synthesizers' own input spectra, and an untrained sparse random projection of the full input), plus absolute band power (Major 1).
\item All four synthesizers retrained on a second simultaneous EEG-fMRI dataset (XP1, OpenNeuro ds002336) and applied to all cohorts (Major 4).
\item Artefact handling: re-harmonisation of every cohort with ICA (ocular components), a band-power view without 30--45\,Hz, artefact measures compared between groups, and a rest-only crop in Parkinson's disease (Major 6).
\item Covariates (age, sex, medication dose, recording hour and year) regressed out within training folds, and recording metadata tests reported (Major 9).
\item Statistics: 1,000 permutations for EEG views and 500 for synthesized views, exact bootstrap intervals at Bonferroni levels, bootstrap $p$-values with Holm correction in stated families (Major 7).
\item A full fidelity table (Table~7), a figure of synthesized volumes (Fig.~5) and corrected fidelity statements (Major 5).
\item Rewritten abstract, Introduction, Results, Discussion and Conclusions; new Fig.~1 (study design); Additional file~2 with supplementary tables.
\end{enumerate}
While re-running the analyses we also found and corrected four issues of our own, which we report for transparency: (i) the CNN-TAG used for the clinical cohorts had been selected after its first training epoch (validation SSIM did not improve later); this is now stated and a 50-epoch model is evaluated as a sensitivity analysis; (ii) the previous Spec2VolCAMU-Net fidelity numbers came from a summary written before its last cross-validation folds had finished (corrected in Table~7); (iii) each label permutation re-ran the nested procedure with two repeats rather than twenty, now stated (this makes the test conservative); (iv) the Parkinson's disease rest file contains a set-up period and an auditory oddball task after the eyes-closed and eyes-open spans, so the fixed 0--400\,s crop was not pure rest; a rest-only crop is now analysed.

\section*{Major comments}

\rc{1. The band-power baseline is not information-matched to the synthesized view \ldots\ I would ask for at least one information-matched EEG baseline.}
\resp We agree, and this was the most consequential change. We added (Section~3.7, Tables~2 and~4, Fig.~4): (a) the temporal mean and SD over pseudo-TRs of the synthesizers' own input spectra (57 channels $\times$ 97 bins), PCA to 20 components exactly as for the synthesized view; (b) the same summary of an untrained sparse random projection (8,192 dimensions) of each pseudo-TR's complete $57\times97\times20$ input; and (c) absolute band power. In Parkinson's disease the information-matched view was the best single view (AUC $0.825$, $p=0.001$, against $0.786$ for band power), and the random projection did equally well ($0.823$). Added to the information-matched view, no trained synthesizer improved classification in any cohort; gains of synthesized fMRI over band power (e.g.\ $+0.035$ for Spec2VolCAMU-Net and $+0.038$ for NT-ViT in Parkinson's disease) were matched by adding the information-matched view to band power ($+0.030$). The reviewer's distinction between ``synthesis adds nothing beyond EEG'' and ``this EEG summary is weak'' is therefore now tested directly, and the answer is the former for the trained synthesizers.

\rc{2. The abstract and conclusions understate a significant result \ldots\ The untrained Semmelweis result should be stated in the abstract and discussed.}
\resp Agreed. The result persists in the revised analysis and now survives Holm correction against both baselines: the untrained Spec2VolCAMU-Net in Semmelweis reaches AUC $0.740$ ($p=0.006$) and adds $+0.165$ [$+0.089$, $+0.248$] over band power and $+0.169$ [$+0.073$, $+0.265$] over the information-matched view; the NODDI-trained network adds $+0.004$. It is stated in the abstract, Results (Section~4.3) and Conclusions, and discussed as an architecture effect in a new Discussion paragraph (``Representation, not information''), as the reviewer suggests. We also tested its robustness, and it did not hold: it shrank to $0.62$ when age, sex and recording hour were regressed out (Semmelweis patients were recorded earlier in the day, $p<0.001$), it disappeared when ocular components were removed by ICA ($0.57$, $p=0.13$), and neither the NODDI- nor the XP1-trained Spec2VolCAMU-Net reproduced it. We now report it as most likely reflecting ocular artefacts and recording conditions that differ between groups in this cohort; the reviewer's question therefore led to a clearer answer than either of us expected.

\rc{3. The conceptual framing of ``Why'' needs sharpening \ldots\ The proposed value of synthesis in small-sample diagnosis is not new information but a better representation or prior.}
\resp We agree and have reframed the hypothesis accordingly in the Introduction (second paragraph) and the Discussion (``Representation, not information''). The controls now separate three possible sources of a gain: a richer EEG summary (information-matched views), the inductive bias of the architecture (untrained networks) and the learned mapping (trained versus untrained). Every gain we observed came from the first two.

\rc{4. The negative result may follow from synthesizer quality and domain shift \ldots\ Ideally, the authors would retrain on at least one further simultaneous EEG-fMRI dataset \ldots\ If that is infeasible, a clear statement of scope is the minimum.}
\resp We did both. All four synthesizers were retrained on XP1 (OpenNeuro ds002336; ten healthy adults, 3\,T, TR 2\,s, near whole-brain coverage, the same 64-channel cap layout as NODDI) with the NODDI training protocol, and applied to all cohorts (Section~3.5, Table~5, Table~S2). We chose XP1 over XP2 (ds002338) because XP2 covers only a 16-slice slab around motor cortex at TR 1\,s. As in NODDI, XP1 required care: the number of raw scanner triggers differed from the number of volumes in 15 of 35 runs, so we used only the 19 runs with exactly one trigger per volume, matched the corrected EEG to the raw triggers through the task markers, and validated the match on all markers. XP1-trained synthesizers reached AUC $0.34$--$0.68$ across cohorts, none improved on the information-matched view, and their fidelity on held-out XP1 subjects was again at the level of an EEG-free template (SSIM $0.662$--$0.699$ against $0.684$). A new ``Scope'' paragraph in the Discussion states that the conclusions concern current volumetric synthesizers trained on small healthy simultaneous datasets and applied across domains, not EEG-to-fMRI synthesis as a concept, and the abstract's conclusion is restricted to ``current synthesized fMRI''.

\rc{5. The fidelity analysis is under-reported \ldots\ Please add a table \ldots\ Please also add a figure of example volumes. \ldots\ The template SSIM is 0.676 in one sentence and 0.675 in the next \ldots}
\resp Added Table~7 (SSIM, PSNR, temporal correlation and shift null for trigger-aligned EEG, file-start EEG and ECG-only input, with per-subject variability and selected epochs; RMSE in Table~S4), Fig.~5 (real volume, template and each synthesizer for a held-out subject, voxelwise temporal SD maps and brain-mean time courses), and the fidelity of the final synthesizers and of the XP1-trained synthesizers on held-out subjects (Table~S4). The two template values are both correct but come from different pipelines: $0.676$ under E2fNet's own data pipeline (EEG paired from the file start) and $0.675$ under our trigger-aligned pipeline; both are now labelled. The reviewer is right that Spec2VolCAMU-Net's published SSIM ($0.693$) is above the template; the sentence now reads that the template is above the published values of E2fNet and NT-ViT and slightly below that of Spec2VolCAMU-Net, and that published values come from each study's own split and preprocessing and are indicative only. In our own leave-one-subject-out evaluation, Spec2VolCAMU-Net ties the template on SSIM ($0.679$ against $0.675$, $p=0.30$) and is worse on RMSE ($p=0.0002$); the other synthesizers are below the template on both.

\rc{6. Artefact handling is not described \ldots\ Please describe any artefact rejection, or add a sensitivity analysis \ldots\ Please also state how the eyes-closed and eyes-open spans in the 400\,s PD crop are ordered.}
\resp The primary analysis applied no artefact rejection, which is now stated (Section~3.3). We added: (a) re-harmonisation of every recording with ICA, removing EOG-correlated components (using only EOG channels present in every recording of a cohort, so that the cleaning cannot differ by group through the channel set), with all views, synthesized ones included, recomputed; (b) band power without the 30--45\,Hz band (the standard ICA muscle criterion flags most components of 45\,Hz low-passed data and could not be used); (c) artefact measures compared between groups (Table~S3); and (d) in Parkinson's disease, a rest-only crop. The between-group comparisons were informative: in Parkinson's disease controls had more ocular activity than patients (EOG power, $p=0.004$), consistent with reduced blinking in Parkinson's disease, whereas accelerometer power did not differ ($p=0.28$); Semmelweis patients had more high-amplitude segments ($p=0.003$); the number of ocular components removed did not differ between groups in any cohort. After ICA, Parkinson's disease was unchanged (band power $0.81$, information-matched $0.80$), band power fell in FEP1 ($0.62\to0.54$) and FEP2 ($0.65\to0.58$), and the untrained Spec2VolCAMU-Net gain in Semmelweis disappeared ($0.57$, $p=0.13$); no synthesized view gave a gain that survived Holm correction in any cohort (smallest adjusted $p=0.42$ over band power and $0.08$ over the information-matched view, 40 comparisons each); the only difference that survived among the trained synthesizers was a loss (Spec2VolCAMU-Net in FEP2, $-0.156$ relative to the information-matched view). Ordering: in every Parkinson's disease recording the file contains a set-up period, one minute of eyes closed, one minute of eyes open and then an auditory oddball task; the onsets of the rest spans and of the task did not differ between groups ($p=0.54$, $p=0.31$), but the 0--400\,s crop included the set-up period and the start of the task. With a rest-only crop, band power still separated the groups ($0.81$; $0.82$ after ICA), so the Parkinson's disease discrimination is not carried by the task, the set-up period or ocular artefacts. In this crop the trained CNN-TAG view reached $0.81$, above its untrained version ($0.57$), the only analysis in which training clearly beat random weights, but it added nothing to band power ($+0.013$, 95\% CI $-0.018$ to $+0.045$) and no comparison in this analysis survived Holm correction.

\rc{7. Some statistical inference choices need justification or strengthening.}
\resp \textit{Multiple comparisons.} The added-value comparisons now use a paired stratified bootstrap with 20,000 resamples, giving percentile intervals at 95\% and at the Bonferroni levels for 20 and 40 comparisons, and a two-sided bootstrap $p$-value, Holm-adjusted within explicitly stated families: the 20 trained-synthesizer comparisons and the 40 comparisons including untrained synthesizers, separately for each baseline (Section~3.7, Table~4, Table~S1). We did not run a permutation test on $\Delta$AUC, because each permutation of the synthesized block would re-run two nested cross-validations; the bootstrap answers the same question without approximation. \textit{Permutation counts.} 1,000 permutations for the EEG views and 500 for the synthesized views (minimum $p\approx0.001$ and $0.002$). The contrastive classifier retrains a network for every permutation, so it keeps 20 permutations, and its $p$-values are now explicitly described as exploratory. Each permutation re-runs the complete nested procedure with two repeats, which widens the null relative to the 20-repeat estimate and makes the test conservative; this is now stated. \textit{Confidence intervals.} Subjects are resampled with their repeat-averaged out-of-fold probabilities (now stated). Principal components are now computed exactly through the subject Gram matrix rather than with a randomised solver; AUCs changed by at most $0.004$.

\rc{8. CNN-TAG in FEP2 needs explanation.}
\resp Section~4.5 reports a diagnostic. None of the first 20 principal components of the CNN-TAG view is associated with sampling rate, recording duration, age, sex or medication dose ($|\rho|\leq0.20$), and no subject carries more than 26\% of any of the first five. In every outer fold the inner cross-validation AUC is below 0.5 ($0.35$--$0.44$) while the training AUC is $0.71$--$0.83$: the signature of anti-learning in features without signal, where a near-zero full-sample group difference forces the difference in a training fold to be opposite to that of its held-out fold. The other synthesized views of FEP2 do not show it.

\rc{9. Clinical confounds should be analysed, not only listed.}
\resp Metadata tests are reported in Table~S3 (FEP1 and Parkinson's disease: no differences; FEP2: longer raw recordings in patients, removed by the fixed crop; Semmelweis: earlier recording hour in patients, $p<0.001$; Fribourg: recording year, all nine 2016 participants were patients, and auxiliary channel layout). Covariates (age, sex; chlorpromazine-equivalent dose in FEP1 and FEP2; recording hour in Semmelweis and Fribourg; recording year in Fribourg) were regressed out of every feature within training folds (Section~3.7, Table~5, Table~S2). Parkinson's disease band power and the information-matched view were unchanged ($0.80$, $0.83$); FEP1 band power fell from $0.62$ to $0.54$; the untrained Spec2VolCAMU-Net in Semmelweis fell from $0.74$ to $0.62$.

\section*{Minor comments}

\rc{Section 3.4 formalism.} Equation~1 is now the binary cross-entropy with a logistic function and an explicit $l_2$ penalty, the set-union notation is removed (the synthesized view is written as the composition $v_2 = F\circ v_1$), and the STFT-view paragraph is rewritten.

\rc{Contrastive loss (Eq.~3).} The absolute deviation $|D-m|$ is the original method's choice and is now justified (Section~3.6): with the layer normalisation around $\pi/2$, pulling mismatched pairs to the margin keeps both classes inside one non-periodic interval of the cosine, whereas a hinge lets distances grow past $\pi$, where the cosine wraps and the classes overlap again after projection. The notation is simplified.

\rc{Figure 3.} The caption's last sentence is completed.

\rc{Borderline interval.} Table~S1 reports all intervals to three decimals; in the revised analysis the Fribourg E2fNet interval over band power is [$+0.004$, $+0.160$] for the NODDI-trained model, and no Fribourg comparison survives correction.

\rc{PCA to 20 components.} Sensitivity to 5, 10 and 40 components and to block-pooled volumes is in Table~5 and Table~S2. Single synthesized views varied by up to $0.4$ AUC across settings (e.g.\ untrained E2fNet in Parkinson's disease, $0.34$ with 5 components and $0.77$ with 40), which we now report as a limitation; no setting brought a trained synthesizer above the information-matched view.

\rc{Pooled FEP result.} Now commented (Section~4.1): pooling lowered the AUC, consistent with heterogeneity between the two samples, which were also recorded mostly at different sampling rates (1 and 3\,kHz).

\rc{Abbreviations and casing.} View names are now used consistently (band-power view, STFT view, synthesized-fMRI view) and CNN-TAG is defined at first use in the main text and in the list of abbreviations.

\rc{Placeholders.} \verify{All placeholders (affiliations, ethics approval IDs, repository, licence, funding, author contributions, competing interests) are completed by the authors before resubmission, and every reference marked for verification has been checked against its source.}

\rc{References.} The Bricman et al.\ entry is now a software citation with its repository URL; entries without volume or page details have been completed \verify{complete after checking}.

\rc{Author guidelines.} The abstract is structured (Background, Methods, Results, Conclusions) and has 247 words; keywords, abbreviations and declarations follow the journal's guidelines.

\end{document}
